\documentclass[11pt]{article}
\usepackage[T1]{fontenc}
\usepackage{lmodern,amsmath,amssymb,amsthm}
\usepackage[margin=1in]{geometry}
\usepackage[hidelinks]{hyperref}
\newtheorem{theorem}{Theorem}
\newtheorem{corollary}[theorem]{Corollary}
\newcommand{\Sha}{\mathrm{Sha}}
\newcommand{\Sel}{\operatorname{Sel}}
\newcommand{\res}{\operatorname{res}}
\newcommand{\cor}{\operatorname{cor}}
\newcommand{\rk}{\operatorname{rank}}
\title{Restriction and Doubling of Two-Primary Tate--Shafarevich Groups in a Quadratic Extension}
\author{\small \href{mailto:admin@cosmosgeek.org}{admin@cosmosgeek.org}\qquad \href{mailto:b321014@iiit-bh.ac.in}{b321014@iiit-bh.ac.in}}
\date{Working paper}
\begin{document}
\maketitle
\begin{abstract}
Consider an elliptic curve and a quadratic twist having identical local
two-Selmer conditions. Suppose the twist's Cassels--Tate pairing and its
sum with the original pairing are both nondegenerate on the common
two-Selmer group. If the splitting field of the twist has no rational
two-torsion on the curve, its Mordell--Weil rank stays fixed and the image
of the original two-primary Tate--Shafarevich group under restriction
is exactly twice the new group. In the finite case, each nontrivial
cyclic elementary divisor doubles in order. For full $S_3$ action on two-torsion and even two-Selmer dimension at least four, Morgan's pairing prescription supplies an infinite tower of such steps. A complementary kernel
for the twist recovers the original group of rational points modulo two.
\end{abstract}

\section{The quadratic step}
Let $K$ be a number field, $D/K$ a quadratic extension, and
$\chi$ its character. Let $E/K$ be an elliptic curve and $E^\chi$
the quadratic twist. Identify $E^\chi[2]$ with $E[2]$. Over $D$ we fix
a twisting isomorphism $E^\chi_D\simeq E_D$ whose conjugate is its negative.
For a place $v$ of $K$, the local conditions below are the Kummer images
of $E(K_v)/2E(K_v)$ and $E^\chi(K_v)/2E^\chi(K_v)$ in $H^1(K_v,E[2])$.
For a field $F$ write $\Sha(E/F)$ for the kernel of localization on
$H^1(F,E)$. A pairing on a two-Selmer group means the pullback of the
Cassels--Tate pairing on $\Sha[2]$, with values in
$\tfrac12\mathbb Z/\mathbb Z$.

\begin{theorem}\label{main}
Assume $E(D)[2]=0$ and that the two local conditions coincide at every
place of $K$. Write $S$ for the common global two-Selmer group and
$C,C_\chi$ for the two pairings on $S$. Assume that $C_\chi$ and
$C+C_\chi$ are both nondegenerate. Put
\[
 A=\Sha(E/K)[2^\infty],\quad B=\Sha(E/D)[2^\infty],\quad
 T=\Sha(E^\chi/K)[2^\infty].
\]
Then
\begin{align*}
 \rk E^\chi(K)&=0,&\rk E(D)&=\rk E(K),\\
 \res:\Sel_2(E/K)&\xrightarrow{\sim}\Sel_2(E/D),&
 B[2]&=\res A[2].
\end{align*}
Furthermore $\res:A\to B$ is injective and
\[
                         \res A=2B.
\]
These conclusions hold for groups with nonzero divisible subgroups as well.
\end{theorem}

\begin{proof}
The kernel of $S\to T[2]$ is $E^\chi(K)/2E^\chi(K)$. It lies in the
radical of $C_\chi$, so it is zero. Since $E^\chi(K)[2]=0$,
finite generation gives $\rk E^\chi(K)=0$. If $T$ had an element of
order four, its double would be a nonzero element of $T[2]$ orthogonal
to all of $T[2]$ by bilinearity. Hence $T=T[2]$. The decomposition of
$E(D)\otimes\mathbb Q$ into the eigenspaces for
$\sigma\in\operatorname{Gal}(D/K)$ now gives $\rk E(D)=\rk E(K)$.

Kramer's theorem identifies the radical of the sum pairing on the
intersection of the two Selmer groups with $\cor\Sel_2(E/D)$
\cite[Thm.~2]{Kramer}, \cite[Thm.~8.4 and Rem.~8.5]{MS}.
In particular, corestriction is zero. For completeness, the
augmentation sequence of the group ring of $\operatorname{Gal}(D/K)$
over $\mathbb F_2$, tensored with $E[2]$, has maps inducing restriction
and corestriction under Shapiro's isomorphism. Its cohomology sequence
is exact at $H^1(D,E[2])$. The inverse image under local restriction
of the Kummer condition over $D_w$ is the sum of the two local
conditions over $K_v$ \cite[Lemmas~2--3]{Kramer};
see also \cite[Rem.~4.4 and Lem.~4.6]{MP}. Under our hypotheses this
sum equals the condition defining $S$. Consequently every class of
$\Sel_2(E/D)$ comes from $S$. Inflation--restriction gives injectivity
of $\res$ on $H^1(K,E[2])$ because $E(D)[2]=0$.

There is an injection $E(K)/2E(K)\to E(D)/2E(D)$. Indeed, for $P=2Q$
with $P\in E(K)$, $Q\in E(D)$, one has $2(\sigma Q-Q)=0$, so $Q\in E(K)$.
Both groups have the same dimension over $\mathbb F_2$ by rank stability
and the two-torsion hypothesis. Comparing the Kummer sequences therefore
identifies $A[2]$ with $B[2]$ by restriction.

To see injectivity on $A$, note that $E(D)_{\rm tors}$ has odd order,
so its positive-degree cohomology for $\operatorname{Gal}(D/K)$ vanishes.
The action on $E(D)/E(D)_{\rm tors}$ is trivial by rank stability,
and $H^1$ of a finite group acting trivially on a free abelian group
is zero. Thus $H^1(D/K,E(D))=0$, and inflation--restriction gives
the asserted injectivity on $A$.

For $b\in B$, ordinary and twisted corestriction give respectively
\[
 (1+\sigma)b=\res\cor b,\qquad
 (1-\sigma)b=\res_\chi\cor_\chi b.
\]
The first lies in $\res A$ and the second lies in $B[2]$ because
$\cor_\chi b\in T$ has exponent two. Since $B[2]=\res A[2]$, adding
the two expressions proves $2B\subseteq\res A$.

Conversely, for $y\in B[2]$ one has $\cor y=0$ because $y=\res x$ for
some $x\in A[2]$. Adjointness of restriction and corestriction under
the pairing \cite[Cor.~4.9]{MS} makes $\res A$ orthogonal to $B[2]$.
We justify $B[2]^\perp=2B$. The Cassels--Tate pairing induces a perfect
pairing on the finite group $B/B_{\rm div}$: the kernel on $B$ is the
maximal divisible subgroup and $B$ has finite corank and finite $B[2]$.
The image of $B[2]$ in that quotient is the full subgroup killed by
two because multiplication by two on $B_{\rm div}$ is surjective. For
a finite group with a perfect pairing, the annihilator of this subgroup
is the image of multiplication by two, by duality. Its preimage in $B$
is $2B$, as $B_{\rm div}\subseteq2B$. This proves the reverse inclusion.
\end{proof}

\section{Consequences for a given extension or tower}
\begin{corollary}
If $A=\bigoplus_{j=1}^{s}\mathbb Z/2^{e_j}\mathbb Z$ is finite with
$e_j\ge1$, then
\[
 B\simeq\bigoplus_{j=1}^{s}\mathbb Z/2^{e_j+1}\mathbb Z.
\]
In a tower of $h$ successive quadratic extensions each satisfying
Theorem~\ref{main}, the exponents become $e_j+h$.
\end{corollary}
\begin{proof}
$2B\simeq A$ and $B[2]\simeq A[2]$ are finite, so $B$ is finite.
The elementary divisors of $2B$ determine those of $B$ larger than two.
Their number already accounts for all $s$ dimensions of $B[2]$.
Iterate the one-step statement.
\end{proof}

\begin{corollary}
There is a natural isomorphism, using the common two-torsion and the
fixed twisting identification over $D$,
\[
 \ker\bigl(T[2]\xrightarrow{\res_\chi}B[2]\bigr)
                  \simeq E(K)/2E(K).
\]
\end{corollary}
\begin{proof}
The first part of the proof of Theorem~\ref{main} identifies
$T[2]$ with the whole of $S$. After restriction to $D$, passage to
$B[2]$ is the projection from $\Sel_2(E/D)$ onto $B[2]$. Its kernel
is $E(D)/2E(D)$, already identified with $E(K)/2E(K)$.
\end{proof}

\section{Prescription of pairings}
The hypotheses specify the pair of local structures and both global
pairings for the given extension. Their linear-algebraic compatibility
can be checked without choosing a twist. In an even-dimensional
$\mathbb F_2$-space $S$ of dimension at least four, every alternating
form $C$ admits an alternating form $P$ for which $P$ and $P+C$ are
nondegenerate. If $\operatorname{rank} C=2h\ge4$, identify a symplectic
complement to its radical with $\mathbb F_{2^h}^2$ equipped with
$\operatorname{Tr}(xy'-x'y)$. For $a\notin\{0,1\}$ the form with
scalar $a$ inside the trace and its sum with $C$ are nondegenerate
there. Give the even-dimensional radical any nondegenerate alternating
form. If $C$ has rank two, pair a symplectic plane with two independent
radical directions by the off-diagonal identity matrix, and use a
symplectic form on the remaining radical. The off-diagonal block keeps
both forms invertible, since its bottom-right block is zero. For $C=0$
one nondegenerate alternating form suffices.

\subsection{Existence of a rank-stable tower}
\begin{theorem}\label{tower}
Let $E/K$ be an elliptic curve over a number field such that
$\operatorname{Gal}(K(E[2])/K)\simeq S_3$. Suppose
$n=\dim_{\mathbb F_2}\Sel_2(E/K)$ is even and $n\ge4$.
Given a finite set of places of $K$, there exists an infinite tower
$K=K_0\subset K_1\subset\cdots$ with quadratic steps, completely split
above the given places and all archimedean places. At every step
\begin{align*}
 \rk E(K_{i+1})&=\rk E(K_i),&
 \res:\Sel_2(E/K_i)&\xrightarrow{\sim}\Sel_2(E/K_{i+1}),\\
 \res\Sha(E/K_i)[2^\infty]&=2\Sha(E/K_{i+1})[2^\infty].
\end{align*}
The restriction on the two-primary Tate--Shafarevich groups is
injective. If the initial group is finite with elementary divisors
$2^{e_1},\ldots,2^{e_s}$, its elementary divisors at $K_i$ are
$2^{e_1+i},\ldots,2^{e_s+i}$.
\end{theorem}

\begin{proof}
We use the pairing prescription in Morgan \cite[Thm.~5.8]{Morgan}
with its local choices. Section~1.1 discusses higher-dimensional
Kummer varieties. Section~2 sets up the subsequent results for a
principally polarized abelian variety over a number field, and
Section~5.2 specifies parts (A) and (B) of Assumption~1.3.
For the natural $S_3$-module $V=E[2]$, the action permutes the three
nonzero vectors transitively, so $V$ is simple. Its centralizing
endomorphism algebra is $\mathbb F_2$: the invertible endomorphisms
commuting with $\operatorname{GL}_2(\mathbb F_2)=S_3$ form its trivial
center, and every nonzero equivariant endomorphism is invertible.
A three-cycle has no nonzero fixed vector. These verify (A) and (B).
The polarization class $\mathfrak c$ vanishes since $\mathcal O_E(O)$
is a rational symmetric line bundle inducing the polarization.
Inflation--restriction for the normal $C_3$ in $S_3$ gives
$H^1(S_3,V)=0$ as $V^{C_3}=0$ and $|C_3|$ is odd.
Thus the intersection hypothesis of Theorem~5.8 holds. By
Definition~5.5 and Lemma~5.7 there, every alternating pairing on $S$
is admissible: it has the same diagonal and the same value on
$(\mathfrak c,\mathfrak c)$ as the original pairing.

Choose $P$ by the preceding linear algebra with $P$ and $P+C$
nondegenerate. In the proof of Theorem~5.8 enlarge Morgan's set
$\Sigma_0$ to include the prescribed places, those dividing
$2N_E\infty$, and the ramification of the governing extension.
Corollary~4.19 and Chebotarev supply a fresh prime
$\mathfrak p\notin\Sigma_0$ with fixed-point-free Frobenius on $V$
and the required pairing switches. The character obtained there
from Proposition~5.2 is trivial at $\Sigma_0$, ramified at
$\mathfrak p$, and unramified except possibly at $\mathfrak p$ and
other good primes of the same fixed-point-free Frobenius type.
Proposition~3.3 gives $C_\chi=P$. At each newly ramified place
the ambient $H^1(K_v,V)$ is zero by Lemma~2.7(i).
At the remaining good odd unramified places both Kummer conditions
are the unramified subgroup by Lemma~2.7(ii). At $\Sigma_0$ the local
twist is trivial. This gives the equality of \emph{every} local
condition used in Theorem~\ref{main}, as in the proof of Lemma~2.8.

Write $D$ for the quadratic field defined by $\chi$. The good odd
prime $\mathfrak p$ is unramified in $K(E[2])$ and ramified in $D$.
Therefore $D\cap K(E[2])=K$. The two-division Galois group over $D$
remains $S_3$, in particular $E(D)[2]=0$. Theorem~\ref{main} applies.
Its Selmer isomorphism preserves $n$, so the same construction over
$D$ gives the next step, with all places above the initially specified
ones included in the new finite set. Induction proves the tower and
the previous corollary gives the elementary divisors.
\end{proof}

Each quadratic step is supplied over its current number field.
The tower is totally real when $K$ is totally real. The theorem
concerns arithmetic rank and the two-primary groups over that tower;
the central order over the original field is a separate question.

\begin{thebibliography}{9}
\bibitem{Kramer} K.~Kramer, \emph{Arithmetic of elliptic curves upon quadratic
extension}, Trans.\ Amer.\ Math.\ Soc.\ \textbf{264} (1981), 121--135.
\href{https://doi.org/10.1090/S0002-9947-1981-0597871-8}{Published article}.
\bibitem{MP} A.~Morgan and R.~Paterson, \emph{On $2$-Selmer groups of twists
after quadratic extension}, J.\ Lond.\ Math.\ Soc.
\href{https://doi.org/10.1112/jlms.12533}{Published article}.
\bibitem{MS} A.~Morgan and A.~Smith, \emph{Field change for the Cassels--Tate
pairing and applications to class groups}, Res.\ Number Theory (2024).
\href{https://doi.org/10.1007/s40993-024-00545-2}{Published article}.
\bibitem{Morgan} A.~Morgan, \emph{Hasse principle for Kummer varieties in the
case of generic $2$-torsion}, Proc.\ Lond.\ Math.\ Soc.
\href{https://doi.org/10.1112/plms.70066}{Published article}.
\end{thebibliography}
\end{document}
